\section{引言}
\label{sec:intro}
\textcite{james_does_2026} 认为使用 \LaTeX 可能提高经济学家对论文研究质量的主观评价。
古语有云：“工欲善其事，必先利其器。\footnote{出自《论语·卫灵公》，其原文为：子贡问为仁。子曰：“工欲善其事，必先利其器。居是邦也，事其大夫之贤者，友其士之仁者。”意思是：工匠要把活干好，必须先让工具锋利；而一个人要在邦国中行仁、干一番事业，就应当去侍奉大夫中的贤能者、结交士人中的仁德之士。}”
一个好的科研工具可以帮助研究者解决中不必要的麻烦，这无疑会提高科研的效率和质量。

\LaTeX 这种为科研排版而生的工具相比于 Microsoft Word （下简称为“MS Word” 或 “Word”）在排版中更胜一筹，且我们不难发现它在给定模版文件的情况下可以轻松地将排版切换至其他期刊。
同时，\LaTeX 这种“源代码-编译-PDF产物”的模式相比于 MS Word 更适合协作和版本追踪。\footnote{更多协作与版本追踪的内容详见第~\ref{sec:github}~节和第~\ref{sec:col}~节。}
此外，它在保持规范化的条件下有更高的灵活性和自由度，表~\ref{tab:word-latex} 对 \LaTeX 和 Word 进行了对比。

\begin{table}[htbp]
    \centering
    \caption{\LaTeX{} 与 Word 在论文写作中的对比}
    \label{tab:word-latex}
    \zihao{5}
    \setlength{\tabcolsep}{3pt}
    \begin{tabular}{@{}p{0.20\linewidth}p{0.38\linewidth}p{0.38\linewidth}@{}}
        \toprule
        \textbf{维度} & \textbf{Microsoft Word} & \textbf{\LaTeX} \\
        \midrule
        \textbf{排版哲学}
          & \textbf{WYSIWYG}（所见即所得），你看到的就是最终打印出来的效果。
          & \textbf{WYSIWYM}（所见即所得意），你只写内容和逻辑结构，由引擎自动排版。 \\
        \addlinespace[2pt]
        \textbf{学习曲线}
          & \textbf{平缓}，上手极快，打开即写，对普通用户极其友好。
          & \textbf{陡峭}，需要学习标记语言、宏包和编译规则，有一定的门槛。 \\
        \addlinespace[2pt]
        \textbf{学习曲线（AI）}
          & AI 较难直接编辑 Word 文档并保留原有格式。
          & AI 可以直接生成和修改源码，作者几乎不需要学习 \LaTeX{} 语法就能开始写作。 \\
        \addlinespace[2pt]
        \textbf{公式与排版}
          & 输入复杂数学公式和大量图表交叉引用时，容易格式错乱或变卡。
          & \textbf{专业且完美}，极其擅长排版长篇巨著、复杂公式、多级定理与交叉引用。 \\
        \addlinespace[2pt]
        \textbf{协作与审阅}
          & \textbf{强}，支持多人实时在线协作、批注、修订（Word/Office 365）。
          & \textbf{代码式协作}，基于纯文本（如 Overleaf 或 Git 版本控制），类似写代码的分支与合并。 \\
        \addlinespace[2pt]
        \textbf{文件格式}
          & 复杂混合格式。
          & \textbf{纯文本文件}（\texttt{.tex}），轻量。 \\
        \addlinespace[2pt]
        \textbf{版本管理}
          & 不方便做细致的版本 Diff。
          & 天然契合 Git 等版本管理工具。 \\
        \addlinespace[2pt]
        \textbf{适用场景}
          & 日常办公、商业报告、简单论文、简历、图文混排、快速修改。
          & 学术论文、学位论文、书籍，以及含有大量公式或交叉引用的科技文档。 \\
        \bottomrule
    \end{tabular}
    \par\smallskip
    {\footnotesize\raggedright
    \textbf{注：}该对比内容由 Gemini 生成，可以代表作者部分观点\par}
\end{table}


在 AI 的辅助中，文档的存储结构影响模型获取和修改文本的方式。
Word 的 \texttt{.docx} 文件采用 Office Open XML（OOXML）格式，正文和样式信息以 XML 形式存储。
提取正文时，需要将分布在多个 XML 节点中的文本片段按文档顺序拼接。
而 \LaTeX 源文件采用纯文本格式，正文内容与排版命令可在同一文件或不同文件中读取，使得内容与格式配置可以解耦，整体版式也可由文档类统一管理。
这种结构有助于简化文本提取流程，使研究者或者模型能够直接在源文件中定位和修改内容。
此前，我们针对于 AI 的能力写过一篇如何让社科研究者用好 \LaTeX 的研究博文，
感兴趣者可以\href{https://marx-max.aidea-labs.com/blog/2026-10-26-latex-guide-social-scientists/}{点击这里查看}。

尽管 \LaTeX 的优势有这么多，但依旧遗憾的是我们在调查了大量的经济学中文期刊后发现，极少数的期刊提供了 \LaTeX 模版，为了让更多的研究者能够从 \LaTeX 中受益，我们做了该项目并决定开源出来。
该项目所采用的开源协议为 \href{https://github.com/aidea-labs/econ-sbm-template?tab=AGPL-3.0-1-ov-file}{AGPL 3.0 License}\footnote{GNU AFFERO GENERAL PUBLIC LICENSE Version 3（AGPL v3.0 License）允许使用、复制、修改和再分发受其约束的作品。分发修改版时，应按同一许可证授权，并依其条款提供相应源代码。若修改版支持通过网络与用户交互，还须向这些用户提供获取相应源代码的途径。}，该许可证允许任何研究者使用该模版。

该项目的边际贡献主要体现在如下三个方面：
（1）针对部分没有提供 \LaTeX{} 模板的中文经济学期刊，本项目参考其公开的投稿要求制作了相应模板，减少研究者自行调整格式的工作；
（2）本项目参考并整合已有模板，尽可能将论文内容与期刊样式分离，使研究者可以在保留同一份论文内容的基础上切换模板；
（3）本项目同时为社会科学研究者提供了一套需要较少技术的 \LaTeX{} 食用方式，介绍公式、图表和文献引用，以及借助 GitHub 和 AI Agent 进行写作与协作的方法（搭配麦当劳薯条 \twemoji{french fries} 食用更美味，此处无广告）。

本文剩余部分安排如下。
第~\ref{sec:works}~章介绍与本项目潜在相关的工作，包括 AEA 的审阅模式和一些 AI Agent 的基础知识。
第~\ref{sec:use}~章介绍如何使用本项目：首先在第~\ref{sec:github}~节介绍 Git 和 GitHub，随后在第~\ref{sec:how}~节介绍如何使用模板撰写论文，最后在第~\ref{sec:col}~节介绍如何通过 GitHub 与合作者共同完成论文。
第~\ref{sec:conclusion}~章对全文进行总结，并讨论项目的局限与后续工作。
